% ============================================================
%  settings.tex --- COPIA ESTO COMPLETO
% ============================================================
\usepackage[utf8]{inputenc}
\usepackage[T1]{fontenc}
\usepackage[spanish,es-noshorthands]{babel}
\usepackage{amsmath, amsfonts, amssymb, amsthm, mathtools, bm, physics}
\usepackage[a4paper, margin=2.5cm]{geometry}
\usepackage{setspace, parskip, microtype, graphicx, booktabs, array, float, caption, xcolor, tcolorbox, enumitem, titlesec, fancyhdr, hyperref}

\tcbuselibrary{theorems, skins, breakable}
\allowdisplaybreaks
\onehalfspacing

% Comandos especiales para fórmulas de las Partes 1 y 4
\newcommand{\E}{\mathbb{E}_t}
\newcommand{\xiit}{\bm{\xi}_t}
\newcommand{\rs}{r^*}
\newcommand{\ns}{n_S}
\newcommand{\nb}{n_B}
\newcommand{\bs}{\beta_S}
\newcommand{\bb}{\beta_B}
\newcommand{\cas}{CA_S}
\newcommand{\cab}{CA_B}

% Encabezados
\pagestyle{fancy}
\fancyhf{}
\rhead{\textit{Macroeconomía Intermedia}}
\lhead{\textsc{Problem Set 1}}
\rfoot{\thepage}

% Definición de Cajas (Usadas en Parte 2 y 3)
\newtcolorbox{resultado}{colback=yellow!10!white, colframe=orange!70!black, fonttitle=\bfseries, title=Resultado, breakable}
\newtcolorbox{condicion}[1][Condición]{colback=blue!5!white, colframe=blue!50!black, fonttitle=\bfseries, title=#1, breakable}
\newtcolorbox{interpretacion}{colback=green!5!white, colframe=green!50!black, fonttitle=\bfseries, title=Interpretación, breakable}
\newtcolorbox{enunciado}{colback=gray!8!white, colframe=gray!50!black, fonttitle=\bfseries, title=Pregunta, breakable}